\documentclass[11pt]{article}
\usepackage[margin=1in]{geometry}
\usepackage{parskip}
\usepackage{hyperref}
\usepackage{enumitem}
\usepackage{xcolor}

\newcommand{\point}[1]{\medskip\noindent\textbf{#1}\par}
\newcommand{\quoted}[1]{\begin{quote}\itshape #1\end{quote}}
\newcommand{\response}{\par\smallskip\noindent\textbf{Response.}\space}

\title{Response to the Editor\\[0.3em]\large Submission f265ac80-65ad-483a-bcfc-21d010e52041}
\author{Rolando Rene Rodriguez, Jr.}
\date{}

\begin{document}
\maketitle

We thank the editor for the review. The revision is the rewrite the decision letter called for: the contribution is reframed as an instantiation of the symmetric-stigmergic schema of \emph{Insects and Agents: Extending the Metaphor} (MABS25), the novel terminology is dropped for pheromone and stigmergy vocabulary, the suggested pheromone propagation over the domain's temporal structure is implemented and its effect reported, the manuscript is cut to 50 pages, and a code audit followed by a complete rerun has put every empirical claim on an accurate footing. The six points are answered in turn below, with pointers into the revised manuscript; section numbers refer to it.

\section*{1. The valuable core: stigmergy with reasoning agents}

\quoted{Stigmergy is useful not only with ant-like agents, but with more sophisticated reasoning mechanisms (here, LLM agents). A simple, straightforward comparison of multiple implementations of a single application with different coordination mechanisms would be very attractive. But that is buried in unnecessary noise.}

\response This is now the spine of the paper. Section~4 (Symmetric-Stigmergic Formalization) states the core directly: foundation-model agents are deployed stigmergically, holding no beliefs or intentions about one another and coordinating only through the shared field, yet bringing open-ended reasoning to each individual patch. The single-application comparison is the whole of the empirical Section~6: five coordination mechanisms (pheromone-field, conversation, hierarchical, sequential, random) run on one application, meeting room scheduling, under identical models, prompts, problem instances, and tick budget, so the only variable is the coordination mechanism. The surrounding noise that buried this comparison has been removed: the five-theorem apparatus is reduced to one lemma, the five-tradition related-work bridge to a single lineage, and the repeated foundation-model reciprocity passages to brief statements in the abstract, introduction, and conclusion.

\section*{2. Terminology}

\quoted{Why introduce a term, ``pressure field,'' previously not used in the swarm intelligence community? I don't see what it adds to the existing terminology.}

\response The novel terminology is dropped. The method is now described throughout in established pheromone and stigmergy vocabulary, using the three-operator language of \emph{Agents, Pheromones, and Mean-Field Models} (AAMAS11): aggregate (the per-region pressure moving average), propagate (the one-hop spillover of Section~4), and evaporate (the per-tick decay of the two pheromone stores). The agent and environment primitives use the MABS25 notation directly: $A(t)$, $E(t)$, $es_i$, $aa(t)$, $ea(t)$, $ESU$, $ESD$, $ESP$, $ASU$, $ASD$, $ASP$. Beyond dropping the coinage, the two surviving terms now rest on a principled distinction from the stigmergy literature rather than on invented vocabulary: we reserve \emph{pheromone} for the deposited marker stores and \emph{pressure} for the sematectonic field the agents sense, which is exactly Parunak's own sematectonic versus sign-based distinction (\emph{A Survey of Environments and Mechanisms for Human-Human Stigmergy}, E4MAS 2005). The phrase ``pressure field'' thus survives only as the name of that sensed sematectonic field; its first appearance, in the abstract and the introduction, now carries this justification and the citation, and Section~4 maps it explicitly to the environment state $es_i$, so the term names a quantity in the schema rather than claiming novelty. The title changes accordingly, to \emph{Symmetric Stigmergy with Foundation-Model Agents: Pheromone-Field Coordination for Decentralized Artifact Refinement}.

\section*{3. Stigmergy literacy}

\quoted{The references to serious SW applications of swarm intelligence, beyond ACO, are extremely limited. A review of the past 10 years of AAMAS and JAAMAS would be very useful if the author wants to publish to this community.}

\response Related work (Section~2) is rebuilt as a single lineage that runs from digital pheromones to language-model agents and engages software applications of stigmergy beyond ant colony optimization. The spine adds the engineering-principles account of natural multi-agent systems (Parunak 1997), the agents-and-mean-fields continuum (Parunak 2011), polyagent and hierarchical-task-network stigmergy (Brueckner 2000; Brueckner, Parunak, and Hilscher 2009), ant-like telecommunications routing (Bonabeau et al.\ 1998), field-based middleware (Mamei and Zambonelli's TOTA), biology-inspired distributed-computing design patterns (Babaoglu et al.), the environment-as-first-class roadmap (the E4MAS volume), and the symmetric-stigmergy schema itself (MABS25). It connects these to the recent AAMAS and JAAMAS strand on environment-mediated and stigmergic coordination for language-model agents, including blackboard architectures, implicit signaling, pheromone-following foragers, stigmergic reinforcement learning, and pheromone-map robot swarms.

\section*{4. The demonstration uses the domain's intrinsic structure}

\quoted{Most applications of stigmergy take advantage of intrinsic structure in the domain. This paper does have structure in that regions are two-hour time slots, but the relations among these slots are not used. One might expect some propagation of pheromone among adjacent regions/time slots to drive meeting shifts, etc. The overall structure seems impoverished.}

\response We implemented pheromone propagation over the temporal structure exactly as suggested, then measured it: at the tested instance scale it does not measurably help, and the paper reports that null as a lower bound on the contribution, not as a failure of the operator. An environment adjacency relation links each within-day time block to its immediate temporal neighbors, with no edges across days, instantiating the $ea(t)$ relation of the schema. Each tick the environment runtime computes an effective activation pressure $g_i = p_i + \kappa \sum_{j \in N(i)} p_j$ with $\kappa = 0.5$ from regions' own pressures, so a high-pressure block recruits proposals in its idle temporal neighbors, the meeting-shift mechanism the comment describes. It is the $ESP$ operator of the schema, applied one hop per tick from own pressures; it affects only which regions receive proposals, while reported pressure stays the sum of own pressures, the separability that supports the convergence lemma. We implemented it both to complete the schema instantiation and as the direct response this comment asks for. The ablation includes a \texttt{no\_propagation} arm, and its result is null: that arm (36.7\%) is statistically indistinguishable from the full system (40.0\%), with a pooled marginal effect of $+4.8$ points ($p = 0.55$). This is what local decomposability predicts: at this instance scale, with one-hop propagation and a separable metric, relieving one block rarely shifts work to a distant block, so activation spillover has little cross-region structure to exploit. The regime where propagation must bite, multi-hop accumulated diffusion, region relations beyond a single temporal chain, and high combinatorial complexity at swarm scale, is precisely what the paper delimits as future work; Section~7 states this once and prominently, framing the null and the win it accompanies as one bound.

\section*{5. Test problem scale}

\quoted{I find the test problems extremely small, at least for a swarming solution, one strength of which is dealing with high combinatorial complexity. (But this may be needed to make solution by other techniques feasible.)}

\response The observation is correct, and the revision states it plainly. The instances are deliberately sized so that the hierarchical and conversational baselines remain feasible to run, the parenthetical reason the comment itself offers; the cost is that the evaluation does not exercise the high-combinatorial regime where swarming is most distinctive. This revision adds no larger-problem experiments and claims none. Section~7 (Discussion) lists problem scale among its limitations and frames the regime characterization as a lower bound on where decentralized coordination helps, not a test at swarm scale; scaling to larger instances is named as future work. Marking the boundary serves the reader better than a token large-instance result that would not settle the question.

\section*{6. Length}

\quoted{I don't believe 76 mss pages are needed to make the key point.}

\response The manuscript is cut to 50 pages. The reductions, in order of size: the five-tradition related-work bridge is collapsed to one lineage spine; the five theorems are reduced to a single convergence lemma plus three short remarks, with the corresponding proofs removed; the dedicated foundation-model reciprocity subsection is removed, the point now made briefly in the abstract, introduction, and conclusion; the ant-colony narrative framing of the escalation results is cut, keeping the mechanism and its table; the ``when to choose each approach'' subsection is removed; and the three societal subsubsections are compressed to a single paragraph. The deleted-mechanism material also goes: a fitness-and-confidence subsystem that the code never read has been removed from the implementation and from the paper, along with the stability theorem and ablation text built on it.

\section*{Accurate reporting of results}

The decision letter asked that results be accurately reported, overstated conclusions rewritten, and limitations fully explained. The accurate account involves a change of numbers, set out below.

An earlier draft reported parity between the decentralized method and hierarchical control on a small pilot. The full grid then showed a regime-dependent gap, and a subsequent audit of the experiment code found two defects in the measurement path: a system-prompt confound that advantaged some baselines under a Latin-square assignment, and a pressure-metric mismatch across strategies. Both are fixed. All strategies now report on one authoritative pressure metric, and every result comes from a complete regenerated grid of 1350 trials (5 strategies $\times$ \{1, 2, 4\} agents $\times$ \{easy, medium, hard\} $\times$ 30 trials), with a separate nine-arm ablation. This three-step history (pilot parity, grid gap, audit and regeneration) is stated in the manuscript itself, in the reporting note that opens Section~6.

Overstated conclusions are rewritten throughout. The headline is scoped to the regime where it holds: decentralized coordination outperforms hierarchical control where strict-greedy region selection collapses into the rejection loop (Section~7), and the gap narrows on the smallest instances where that loop does not trigger. Final-pressure ratios are reported with the caveat that they average over unsolved trials and so partly measure how often a baseline fails to start; convergence-speed ratios are reported with the small solved-count caveat; effect sizes are reported per comparison rather than as an inequality; and the easy-difficulty ablation is read with its ceiling-effect caveat. All figures and tables report the complete regenerated results, with confidence intervals throughout.

\bigskip
\noindent The revision addresses each of the six points directly.

\end{document}
